\section{The Majumdar--Ghosh state}\label{sec:majumdar_ghosh}
%% ===================================================================

The Majumdar--Ghosh state is the exact dimer ground state of the spin-$\tfrac12$
antiferromagnetic Heisenberg chain with next-nearest-neighbour coupling at the
Majumdar--Ghosh point $J_2 = J_1 / 2$~\cite{MajumdarGhosh1969,MajumdarGhosh1969II}:
$H=\sum \boldsymbol{S}_i\cdot\boldsymbol{S}_{i+1}
    +\tfrac12\sum\boldsymbol{S}_i\cdot\boldsymbol{S}_{i+2}$.
On an even periodic chain it is the symmetric superposition of the two
nearest-neighbour singlet coverings $(1,2)(3,4)\cdots$ and
$(2,3)(4,5)\cdots(N,1)$, forming the one-dimensional prototype of a
resonating valence-bond state. The review
\cite[Appendix~A, ``The Majumdar-Ghosh model'']{Cirac2021Matrix} writes it as a
matrix product state with a tensor built from the singlet matrix
\begin{align}
    Y & =
    \begin{pmatrix}
        0 & -1 \\
        1 & 0
    \end{pmatrix},
    \label{eq:mg_singlet_matrix}
\end{align}
whose entry $Y_{ab}$ is the coefficient of $\ket{ab}$ in the singlet
$\ket{10}-\ket{01}$. The tensor is a non-normal, non-injective example with a
two-fold degenerate periodic ground space.

\begin{definition}[Majumdar--Ghosh tensor of the review]%
    \label{def:majumdar_ghosh_review_tensor}
    \lean{MPSTensor.majumdarGhoshSingletY, MPSTensor.majumdarGhoshReviewBond,
        MPSTensor.majumdarGhoshReviewTensor}
    \leanok
    \uses{def:mps_tensor}
    Set $d=2$. Let $E_{ab}$ denote the matrix units of $\MN{3}$, and let
    \begin{align}
        B^0 & = E_{02}+E_{20},
            & B^1              & = E_{12}+E_{21}.
        \label{eq:mg_review_bond}
    \end{align}
    The review prints the tensor
    $A=\bigl[\ket0\bigl[(02|+(20|\bigr]+\ket1\bigl[(12|+(21|\bigr]\bigr]\otimes Y$
    \cite[Appendix~A, ``The Majumdar-Ghosh model'']{Cirac2021Matrix}. Read
    literally, with $(ab|$ the bra $E_{ab}$ of the pair of virtual indices and
    $\otimes$ the Kronecker product, it is the tensor
    $\widetilde A^i=B^i\otimes Y\in\MN{6}$ on the bond
    $\C^3\otimes\C^2\cong\C^6$, with $Y$ as in~(\ref{eq:mg_singlet_matrix}).
\end{definition}

\begin{lemma}[Coefficients of the review tensor factor]%
    \label{lem:majumdar_ghosh_review_coeff}
    \lean{MPSTensor.majumdarGhoshReview_coeff}
    \leanok
    \uses{def:majumdar_ghosh_review_tensor, def:eval_word}
    For every word $w$,
    $\tr\bigl(\widetilde A^{w}\bigr)=\tr\bigl(B^{w}\bigr)\,\tr\bigl(Y^{|w|}\bigr)$.
\end{lemma}

\begin{proof}
    \leanok
    By induction on $w$ and the mixed-product rule
    $(B^i\otimes Y)(M\otimes Y^n)=B^iM\otimes Y^{n+1}$, the word evaluation is
    $\widetilde A^{w}=B^{w}\otimes Y^{|w|}$ up to the identification of the
    bond with $\C^6$, which preserves traces. The trace of a Kronecker product
    is the product of the traces.
\end{proof}

\begin{definition}[Pair coverings]\label{def:majumdar_ghosh_pair_covering}
    \lean{MPSTensor.pairProduct, MPSTensor.pairCoveringEven,
        MPSTensor.pairCoveringOdd, MPSTensor.majumdarGhoshSinglet}
    \leanok
    Let $f(a,b)$ be a pair amplitude on two sites. For a chain of $N$ sites
    with configurations $\sigma=(\sigma_0,\ldots,\sigma_{N-1})$ put
    \begin{align}
        P^{\mathrm{even}}_f(\sigma) & = \prod_{k=0}^{\lfloor N/2\rfloor-1}
        f(\sigma_{2k},\sigma_{2k+1}),
                                    &
        P^{\mathrm{odd}}_f(\sigma)  & = \prod_{k=0}^{\lfloor N/2\rfloor-1}
        f(\sigma_{2k+1},\sigma_{(2k+2)\bmod N}).
        \label{eq:mg_pair_covering}
    \end{align}
    For even $N$ these are the coverings $(1,2)(3,4)\cdots(N-1,N)$ and
    $(2,3)(4,5)\cdots(N,1)$ of the periodic chain by the pair state
    $\sum_{ab}f(a,b)\ket{ab}$. The normalized singlet amplitude is
    $s(a,b)=Y_{ab}/\sqrt2$. For a word $w$, the pair product is the product of
    $f$ over the consecutive pairs $(w_0,w_1),(w_2,w_3),\ldots$, and zero if
    $|w|$ is odd.
\end{definition}

\begin{lemma}[Pair-covering form of periodic vectors]%
    \label{lem:mps_pair_covering_form}
    \lean{MPSTensor.evalWord_apply_hub_eq_pairProduct,
        MPSTensor.coeff_cons_eq_pairProduct_add,
        MPSTensor.mpv_eq_pairCovering_of_hub,
        MPSTensor.mpv_eq_zero_of_hub_of_odd, MPSTensor.pairProduct_eq_prod,
        MPSTensor.pairProduct_eq_zero_of_odd,
        MPSTensor.pairProduct_ofFn_eq_pairCoveringEven,
        MPSTensor.pairProduct_rotate_ofFn_eq_pairCoveringOdd}
    \leanok
    \uses{def:mpv, def:majumdar_ghosh_pair_covering}
    Let $K$ be a tensor with a bond level $z$ such that $K^i_{xy}=0$ whenever
    $x\neq z$ and $y\neq z$, and $K^i_{zz}=0$ for every $i$. Put
    $f(a,b)=(K^aK^b)_{zz}$. Then $(K^w)_{zz}$ is the pair product of $w$. For
    even $N>0$,
    \begin{align}
        \ket{V^{(N)}(K)} & = P^{\mathrm{even}}_f+P^{\mathrm{odd}}_f,
        \label{eq:mg_pair_vector}
    \end{align}
    and for odd $N$ the periodic vector vanishes.
\end{lemma}

\begin{proof}
    \leanok
    The row $z$ of $K^aK^b$ is $f(a,b)$ times the $z$-th unit row, since a path
    leaving $z$ must enter a level $x\neq z$ and then return to $z$. Induction
    on the word in steps of two letters gives the entry $(K^w)_{zz}$, and a word
    of length one has $(K^a)_{zz}=0$. For a word $a\,w$,
    \begin{align}
        \tr(K^aM) & = (K^aM)_{zz}+(MK^a)_{zz},
        \qquad M=K^{w},
        \label{eq:mg_pair_trace}
    \end{align}
    because the rows $x\neq z$ of $K^a$ are supported in column $z$ and
    $K^a_{zz}=0$. The two terms are the $(z,z)$ entries for $a\,w$ and for its
    rotation $w\,a$, which are the two pair coverings of the configuration.
    Pair products of words of odd length vanish.
\end{proof}

\begin{definition}[Review tensor with the singlet on the spin levels]%
    \label{def:majumdar_ghosh_bond_singlet_tensor}
    \lean{MPSTensor.majumdarGhoshBondSingletTensor}
    \leanok
    \uses{def:majumdar_ghosh_review_tensor}
    The tensor $C^i=B^i\,(Y\oplus1)\in\MN{3}$, with $B^i$ as
    in~(\ref{eq:mg_review_bond}), places the singlet~(\ref{eq:mg_singlet_matrix})
    on the two spin levels $0,1$ of the bond and carries the empty level $2$
    through. Explicitly $C^0=E_{02}-E_{21}$ and $C^1=E_{12}+E_{20}$.
\end{definition}

\begin{theorem}[The singlet-bond tensor gives the dimer superposition]%
    \label{thm:majumdar_ghosh_bond_singlet_mpv}
    \lean{MPSTensor.majumdarGhoshBondSinglet_mpv_eq_pairCovering}
    \leanok
    \uses{def:majumdar_ghosh_bond_singlet_tensor, lem:mps_pair_covering_form}
    For every even $N>0$, $\ket{V^{(N)}(C)}=P^{\mathrm{even}}_Y+P^{\mathrm{odd}}_Y$,
    the sum of the two nearest-neighbour coverings by the singlet
    $\ket{10}-\ket{01}$. This is the Majumdar--Ghosh superposition of
    \cite[Appendix~A, ``The Majumdar-Ghosh model'']{Cirac2021Matrix}, with the
    tensor read as in \cite{gap:rmp_majumdar_ghosh_tensor_gap}.
\end{theorem}

\begin{proof}
    \leanok
    The level $z=2$ satisfies the hypotheses of
    Lemma~\ref{lem:mps_pair_covering_form} for $C$, and $(C^aC^b)_{22}=Y_{ab}$ by
    direct multiplication.
\end{proof}

The representative used below has the same structure, with the distinguished
level $0$ and the normalized singlet as pair amplitude.

\begin{definition}[Majumdar--Ghosh tensor]\label{def:majumdar_ghosh_tensor}
    \lean{MPSTensor.majumdarGhoshTensor}
    \leanok
    \uses{def:mps_tensor}
    Set $d = 2$ (spin-$\tfrac12$) and $D = 3$. The Majumdar--Ghosh tensor is
    the family $\{A^0, A^1\} \subset \MN{3}$ defined by
    \begin{align}
        A^0 & =
        \begin{pmatrix}
            0               & 1 & 0 \\
            0               & 0 & 0 \\
            \tfrac1{\sqrt2} & 0 & 0
        \end{pmatrix},
        \qquad
        A^1 =
        \begin{pmatrix}
            0                & 0 & 1 \\
            -\tfrac1{\sqrt2} & 0 & 0 \\
            0                & 0 & 0
        \end{pmatrix}.
        \notag
    \end{align}
    The level $0$ is the closed link between two dimers, the levels $1,2$ carry
    the spin of an open singlet partner, and the relative sign
    $-\tfrac1{\sqrt2}$ is the antisymmetry of the singlet.
\end{definition}

\begin{theorem}[The Majumdar--Ghosh tensor gives the dimer superposition]%
    \label{thm:majumdar_ghosh_dimer_superposition}
    \lean{MPSTensor.majumdarGhosh_mpv_eq_pairCovering,
        MPSTensor.majumdarGhosh_mpv_eq_zero_of_odd,
        MPSTensor.majumdarGhoshTensor_mul_apply_zero_zero}
    \leanok
    \uses{def:majumdar_ghosh_tensor, lem:mps_pair_covering_form}
    For every even $N>0$, $\ket{V^{(N)}(A)}=P^{\mathrm{even}}_s+P^{\mathrm{odd}}_s$,
    the sum of the two nearest-neighbour coverings by the normalized singlet
    $(\ket{10}-\ket{01})/\sqrt2$. For odd $N$ the periodic vector vanishes.
\end{theorem}

\begin{proof}
    \leanok
    The level $z=0$ satisfies the hypotheses of
    Lemma~\ref{lem:mps_pair_covering_form} for $A$: rows $1,2$ of $A^0$ and $A^1$
    are supported in column $0$, and $A^i_{00}=0$. Direct multiplication gives
    $(A^aA^b)_{00}=Y_{ab}/\sqrt2$.
\end{proof}

\begin{lemma}[Relation to the singlet-bond tensor]%
    \label{lem:majumdar_ghosh_bond_singlet_scaling}
    \lean{MPSTensor.majumdarGhoshBondSinglet_mpv_eq_smul}
    \leanok
    \uses{def:majumdar_ghosh_tensor, thm:majumdar_ghosh_bond_singlet_mpv,
        thm:majumdar_ghosh_dimer_superposition}
    For every even $N>0$,
    $\ket{V^{(N)}(C)}=\sqrt2^{\,N/2}\,\ket{V^{(N)}(A)}$.
\end{lemma}

\begin{proof}
    \leanok
    Both vectors are sums of the two pair coverings, with amplitudes $Y_{ab}$
    and $Y_{ab}/\sqrt2$; each covering is a product of $N/2$ amplitudes.
\end{proof}

\begin{theorem}[The literal review tensor does not give the dimer superposition]%
    \label{thm:majumdar_ghosh_review_mpv}
    \lean{MPSTensor.majumdarGhoshReview_mpv_eq,
        MPSTensor.majumdarGhoshReview_mpv_ne_smul,
        MPSTensor.majumdarGhoshReviewBond_mul_apply_two_two,
        MPSTensor.majumdarGhoshSingletY_sq,
        MPSTensor.majumdarGhoshSingletY_pow_two_mul}
    \leanok
    \uses{def:majumdar_ghosh_review_tensor, lem:majumdar_ghosh_review_coeff,
        lem:mps_pair_covering_form, thm:majumdar_ghosh_dimer_superposition}
    Let $N=2m>0$ and $\delta(a,b)=\delta_{ab}$. Then
    \begin{align}
        \ket{V^{(N)}(\widetilde A)} & =
        2(-1)^m\bigl(P^{\mathrm{even}}_\delta+P^{\mathrm{odd}}_\delta\bigr),
        \label{eq:mg_review_vector}
    \end{align}
    the sum of the two coverings by the pair state $\ket{00}+\ket{11}$, and
    $\ket{V^{(N)}(\widetilde A)}$ is not a scalar multiple of
    $\ket{V^{(N)}(A)}$.
\end{theorem}

\begin{proof}
    \leanok
    The level $z=2$ satisfies the hypotheses of
    Lemma~\ref{lem:mps_pair_covering_form} for $\{B^0,B^1\}$, with
    $(B^aB^b)_{22}=\delta_{ab}$,
    so Lemma~\ref{lem:mps_pair_covering_form} gives
    $\tr(B^{\sigma})=P^{\mathrm{even}}_\delta(\sigma)+P^{\mathrm{odd}}_\delta(\sigma)$.
    Since $Y^2=-I_2$, $\tr(Y^{2m})=2(-1)^m$, and
    Lemma~\ref{lem:majumdar_ghosh_review_coeff} gives~(\ref{eq:mg_review_vector}).
    On the configuration with all spins $0$ the right side
    of~(\ref{eq:mg_review_vector}) equals $4(-1)^m\neq0$, while
    Theorem~\ref{thm:majumdar_ghosh_dimer_superposition} gives
    $V^{(N)}(A)_{0\cdots0}=0$ because $s(0,0)=0$.
\end{proof}

\begin{theorem}[Majumdar--Ghosh tensor is in left-canonical form]%
    \label{thm:majumdar_ghosh_left_canonical}
    \lean{MPSTensor.majumdarGhosh_left_canonical}
    \leanok
    \uses{def:majumdar_ghosh_tensor}
    The Majumdar--Ghosh tensor is isometric:
    $(A^0)^\dagger A^0 + (A^1)^\dagger A^1 = \Id$.
\end{theorem}

\begin{proof}
    \leanok
    Direct multiplication of the two displayed matrices gives
    $(A^0)^\dagger A^0=\diag(\tfrac12,1,0)$ and
    $(A^1)^\dagger A^1=\diag(\tfrac12,0,1)$.
\end{proof}

\begin{theorem}[Transfer map of the Majumdar--Ghosh tensor]%
    \label{thm:majumdar_ghosh_transfer}
    \lean{MPSTensor.majumdarGhosh_transferMap_one}
    \leanok
    \uses{def:majumdar_ghosh_tensor, def:transfer_map}
    The transfer map of the Majumdar--Ghosh tensor sends the identity to the
    diagonal matrix $\E_A(\Id) = \diag(2, \tfrac12, \tfrac12)$.
\end{theorem}

\begin{proof}
    \leanok
    Substitute the two matrices in
    $\E_A(\Id)=A^0(A^0)^\dagger+A^1(A^1)^\dagger$ and multiply.
\end{proof}

\begin{theorem}[Majumdar--Ghosh tensor is not injective]%
    \label{thm:majumdar_ghosh_not_injective}
    \lean{MPSTensor.majumdarGhosh_not_isInjective}
    \leanok
    \uses{def:majumdar_ghosh_tensor, def:injective}
    The Majumdar--Ghosh tensor is not injective: every element of
    $\spn_{\C}\{A^0, A^1\}$ has vanishing $(2,2)$ entry, so the span is a proper
    subspace of~$\MN{3}$.
\end{theorem}

\begin{proof}
    \leanok
    The matrix unit $E_{22}$ has nonzero $(2,2)$ entry, while every linear
    combination of $A^0$ and $A^1$ has zero $(2,2)$ entry.
\end{proof}

\begin{theorem}[Majumdar--Ghosh tensor is not normal]%
    \label{thm:majumdar_ghosh_not_normal}
    \lean{MPSTensor.majumdarGhosh_not_isNormal}
    \leanok
    \uses{def:majumdar_ghosh_tensor, def:normal}
    The Majumdar--Ghosh tensor is not normal: no blocking length makes the
    products of that length span the full matrix algebra~$\MN{3}$. This
    non-normality is the algebraic signature expected from the two-periodic
    dimer-covering structure.
\end{theorem}

\begin{proof}
    \leanok
    The bond carries a two-periodic grading with complementary support sets
    \begin{align}
         & \{(0,1),(0,2),(1,0),(2,0)\}, \notag \\
         & \{(0,0),(1,1),(1,2),(2,1),(2,2)\}.
        \notag
    \end{align}
    Left multiplication with a generator exchanges these sets. Let $E_{ij}$
    denote the $(i,j)$ matrix
    unit. Then, for every blocking length $N$,
    \begin{align}
        N\text{ odd}
         & \Rightarrow
        (\forall w,\ |w|=N \Rightarrow (A^w)_{22}=0)
        \Rightarrow
        E_{22}\notin\spn_{\C}\{A^w: |w|=N\},
        \notag         \\
        N\text{ even}
         & \Rightarrow
        (\forall w,\ |w|=N \Rightarrow (A^w)_{10}=0)
        \Rightarrow
        E_{10}\notin\spn_{\C}\{A^w: |w|=N\}.
        \notag
    \end{align}
    Hence no fixed-length word span is all of $\MN{3}$.
\end{proof}

\begin{remark}\label{rem:majumdar_ghosh_bond_dimension}
    Read literally, the valence-bond formula of \cite{Cirac2021Matrix}
    specifies the six-dimensional tensor of
    Definition~\ref{def:majumdar_ghosh_review_tensor}, which by
    Theorem~\ref{thm:majumdar_ghosh_review_mpv} gives coverings by
    $\ket{00}+\ket{11}$ rather than by singlets. Placing the singlet on the spin
    levels of a three-level bond
    (Definition~\ref{def:majumdar_ghosh_bond_singlet_tensor}) gives the
    Majumdar--Ghosh superposition; see \cite{gap:rmp_majumdar_ghosh_tensor_gap}.
\end{remark}

\begin{definition}[Majumdar--Ghosh Hamiltonian]\label{def:majumdar_ghosh_hamiltonian}
    \lean{MPSTensor.spinHalfOperator, MPSTensor.spinExchange,
        MPSTensor.majumdarGhoshTerm, MPSTensor.majumdarGhoshHamiltonian}
    \leanok
    Let $S^\alpha=\sigma^\alpha/2$ for $\alpha=x,y,z$ be the spin-$\tfrac12$
    operators, and for two sites $j,k$ of a chain let
    $\mathbf S_j\cdot\mathbf S_k=\sum_\alpha S^\alpha_jS^\alpha_k$. On a periodic
    chain of $N$ sites, with indices modulo $N$,
    \begin{align}
        H_N & = \sum_{i}\mathbf S_i\cdot\mathbf S_{i+1}
        +\tfrac12\sum_{i}\mathbf S_i\cdot\mathbf S_{i+2}
        \label{eq:mg_hamiltonian}
    \end{align}
    \cite[Appendix~A, ``The Majumdar-Ghosh model'']{Cirac2021Matrix}. The
    three-site term is
    $h=\tfrac12(\mathbf S_1\cdot\mathbf S_2+\mathbf S_2\cdot\mathbf S_3
        +\mathbf S_1\cdot\mathbf S_3)$.
\end{definition}

\begin{lemma}[Exchange as a transposition]\label{lem:spin_exchange_swap}
    \lean{MPSTensor.spinExchange_apply, MPSTensor.spinHalfOperator_sum}
    \leanok
    \uses{def:majumdar_ghosh_hamiltonian}
    For distinct sites $j\neq k$,
    $\mathbf S_j\cdot\mathbf S_k=\tfrac12P_{jk}-\tfrac14$, where $P_{jk}$
    exchanges the spins at $j$ and $k$.
\end{lemma}

\begin{proof}
    \leanok
    The completeness relation
    $\sum_\alpha S^\alpha_{pa}S^\alpha_{qb}
        =\tfrac12\delta_{pb}\delta_{qa}-\tfrac14\delta_{pa}\delta_{qb}$ is checked on
    the sixteen index values.
\end{proof}

\begin{theorem}[Three-site ground space of the Majumdar--Ghosh tensor]%
    \label{thm:majumdar_ghosh_local_ground_space}
    \lean{MPSTensor.majumdarGhosh_groundSpace_three_eq_eigenspace}
    \leanok
    \uses{def:majumdar_ghosh_tensor, def:ground_space_parent,
        def:majumdar_ghosh_hamiltonian, lem:spin_exchange_swap}
    The three-site local ground space $\mathcal G_3(A)$ equals the eigenspace
    of $h$ for the eigenvalue $-\tfrac38$.
\end{theorem}

\begin{proof}
    \leanok
    By Lemma~\ref{lem:spin_exchange_swap},
    $h=\tfrac14(P_{12}+P_{23}+P_{13})-\tfrac38$, so the eigenspace consists of
    the vectors $v$ with $v(\sigma P_{12})+v(\sigma P_{23})+v(\sigma P_{13})=0$
    for every $\sigma$, that is,
    \begin{align}
        v_{000} & = v_{111} = 0,
                & v_{001}+v_{010}+v_{100} & = 0,
                & v_{011}+v_{101}+v_{110} & = 0.
        \label{eq:mg_spin_half_conditions}
    \end{align}
    Every $\Gamma_3(X)$ satisfies~(\ref{eq:mg_spin_half_conditions}) by
    expansion. Conversely, a vector satisfying these conditions is
    $\Gamma_3(X)$ for the matrix $X$ with $X_{01}=-2v_{110}$,
    $X_{02}=-2v_{001}$, $X_{10}=\sqrt2\,v_{100}$, $X_{20}=-\sqrt2\,v_{011}$ and
    all other entries zero.
\end{proof}

\begin{lemma}[Twisted periodic vectors satisfy the window constraints]%
    \label{lem:twisted_mpv_chain_ground_space}
    \lean{MPSTensor.twistedMPV_window_mem_groundSpace,
        MPSTensor.twistedMPV_mem_chainGroundSpace}
    \leanok
    \uses{def:chain_ground_space}
    Let $G\in\MN{D}$ and $c\in\C$ satisfy $GA^i=c\,A^iG$ for every $i$. Then
    for $0<N$ and $L\le N$ the vector
    $\sigma\mapsto\tr(A^{\sigma_0}\cdots A^{\sigma_{N-1}}G)$ lies in
    $\mathcal G_{N,L}(A)$.
\end{lemma}

\begin{proof}
    \leanok
    For the window starting at site $i$, split the configuration as $u\,v$ with
    $|u|=i$, so that $v\,u$ begins with the window. Then
    $\tr(A^uA^vG)=\tr(A^vGA^u)=c^{i}\tr(A^vA^uG)$, and $A^{vu}G$ is the window
    word followed by the boundary matrix $c^i A^{\mathrm{out}}G$, where
    $A^{\mathrm{out}}$ is the product over the sites outside the window.
\end{proof}

\begin{theorem}[Each dimer covering lies in the chain ground space]%
    \label{thm:majumdar_ghosh_coverings_chain_ground_space}
    \lean{MPSTensor.majumdarGhosh_pairCoveringEven_mem_chainGroundSpace,
        MPSTensor.majumdarGhosh_pairCoveringOdd_mem_chainGroundSpace}
    \leanok
    \uses{thm:majumdar_ghosh_dimer_superposition, lem:mps_pair_covering_form,
        lem:twisted_mpv_chain_ground_space, lem:mpv_mem_chain_ground_space}
    For even $N>0$ and $L\le N$, both coverings $P^{\mathrm{even}}_s$ and
    $P^{\mathrm{odd}}_s$ lie in $\mathcal G_{N,L}(A)$.
\end{theorem}

\begin{proof}
    \leanok
    The grading $G=\diag(1,-1,-1)$ satisfies $GA^i=-A^iG$. For even $N$,
    $\tr(A^{\sigma}G)=2(A^{\sigma})_{00}-\tr(A^{\sigma})
        =P^{\mathrm{even}}_s(\sigma)-P^{\mathrm{odd}}_s(\sigma)$ by
    Lemma~\ref{lem:mps_pair_covering_form} and
    Theorem~\ref{thm:majumdar_ghosh_dimer_superposition}. The two coverings
    are half the sum and half the difference of this vector and
    $\ket{V^{(N)}(A)}$, both of which lie in $\mathcal G_{N,L}(A)$ by
    Lemma~\ref{lem:twisted_mpv_chain_ground_space}.
\end{proof}

\begin{theorem}[Energy of the dimer coverings]%
    \label{thm:majumdar_ghosh_hamiltonian_eigen}
    \lean{MPSTensor.majumdarGhoshHamiltonian_apply_of_mem_chainGroundSpace,
        MPSTensor.majumdarGhoshHamiltonian_pairCoveringEven,
        MPSTensor.majumdarGhoshHamiltonian_pairCoveringOdd,
        MPSTensor.majumdarGhoshHamiltonian_mpv}
    \leanok
    \uses{def:majumdar_ghosh_hamiltonian, thm:majumdar_ghosh_local_ground_space,
        thm:majumdar_ghosh_coverings_chain_ground_space,
        lem:mpv_mem_chain_ground_space}
    For $N\ge3$, every $\psi\in\mathcal G_{N,3}(A)$ satisfies
    $H_N\psi=-\tfrac{3N}{8}\psi$. In particular, for even $N\ge4$ each of the
    two singlet coverings, and for $N\ge3$ the periodic vector
    $\ket{V^{(N)}(A)}$, is an eigenvector of~(\ref{eq:mg_hamiltonian}) with
    eigenvalue $-\tfrac{3N}{8}$.
\end{theorem}

\begin{proof}
    \leanok
    Each window of three sites $i,i+1,i+2$ of $\psi$ lies in
    $\mathcal G_3(A)$, so by
    Theorem~\ref{thm:majumdar_ghosh_local_ground_space} the three transposed
    copies of $\psi$ in that window sum to zero. Summing over $i$ and shifting
    the index of the middle transposition gives
    $2\sum_iP_{i,i+1}\psi+\sum_iP_{i,i+2}\psi=0$. Lemma~\ref{lem:spin_exchange_swap}
    then gives
    $H_N\psi=\tfrac14\bigl(2\sum_iP_{i,i+1}\psi+\sum_iP_{i,i+2}\psi\bigr)
        -\tfrac{3N}{8}\psi=-\tfrac{3N}{8}\psi$.
\end{proof}

% The lower bound $H_N\ge-3N/8$, which would make these eigenvectors ground
% states, and the dimension count below are not yet formalized.
Because the Majumdar--Ghosh tensor is non-normal
(Theorem~\ref{thm:majumdar_ghosh_not_normal}) with a two-periodic bond grading,
the expected two-fold dimer degeneracy of its periodic ground space is stated
in the following target theorem.

\begin{theorem}[Two-fold dimer degeneracy of the Majumdar--Ghosh ground space]%
    \label{thm:majumdar_ghosh_ground_space}
    \notready
    \uses{def:majumdar_ghosh_tensor, thm:majumdar_ghosh_not_normal, def:ground_space_parent,
        thm:majumdar_ghosh_coverings_chain_ground_space}
    On an even periodic chain of length $N = 2m\ge4$, let
    $H_{N,\mathrm{MG}}^{(3)}$ be the range-three parent Hamiltonian of the
    Majumdar--Ghosh tensor, the periodic sum of translates of the local
    projector with kernel $\mathcal G_3(A_{\mathrm{MG}})$. Then
    $H_{N,\mathrm{MG}}^{(3)}$ has a two-dimensional ground space, spanned by
    the two nearest-neighbour singlet coverings
    $(1,2)(3,4)\cdots(N-1,N)$ and $(2,3)(4,5)\cdots(N,1)$.
    This two-fold degeneracy reflects the spontaneous breaking of the one-site
    translation symmetry of the spin-$\tfrac12$ chain
    $H=\sum \boldsymbol{S}_i\cdot\boldsymbol{S}_{i+1} +\tfrac12\sum\boldsymbol{S}_i\cdot\boldsymbol{S}_{i+2}$
    of \cite{Cirac2021Matrix}.
\end{theorem}

%% ===================================================================
